\uuid{4cyC}
\titre{Congruences}
\auteur{bodin ; croizat ; sacre}
\organisation{exo7}
\licence{CC BY-NC-SA 4.0}
\datecreate{2022-03-08}
\niveau{L1}
\module{Algèbre}
\chapitre{Arithmétique dans Z}
\sousChapitre{Autre}
\difficulte{3}
\theme{Congruences}
\exo7qcmid{hilisit.arithmetique.048}

\contenu{
\begin{question}
\qclassification{103.99}
    Soit $p$ un nombre premier et $x$ un entier.
    Quel(s) énoncé(s) du petit théorème de Fermat sont corrects ?
    \begin{answers}
        \bad{$x^p \equiv p \;[x]$}
        \good{$x^p \equiv x \;[p]$}
        \bad{Si $p$ ne divise pas $x$, alors $x^{p-1} \equiv 0 \;[x]$}
        \bad{Si $p$ ne divise pas $x$, alors $x^{p-1} \equiv 0 \;[p]$}
    \end{answers}
    \begin{explanations}
    Le théorème de Fermat stipule que $x^p \equiv x \;[p]$, et que si $p$ ne divise pas $x$ alors $x^{p-1} \equiv 1 \;[p]$.
    \end{explanations}
\end{question}
}
